% Part: lambda-calculus
% Chapter: introduction
% Section: overview

\documentclass[../../../include/open-logic-section]{subfiles}

\begin{document}

\olfileid{lam}{int}{ovr}
\olsection{جائزہ}

لامبڈا کیلکولس نوں Alonzo Church نے 1930 دی دہائی دے شروع وچ
تعمیری منطق دی بنیاد لئی تیار کیتا سی،قابلِ حساب فنکشوناں دے نمونے
وَجوں \emph{نہیں}۔ پر چھیتی ایہ وکھایا گیا کہ ایہ قابلِ حسابیت دیاں
ہور تعریفاں،مثلاً ٹیورنگ دے ذریعے قابلِ حساب فنکشوناں تے جزوی بازگشتی
فنکشوناں،دے ہم ارز اے۔ایہ گل کہ شروع وچ ایس نتیجے اُتے تھوڑی حیرانی
ہوئی سی، ایس ہم ارزی نوں ہور دلچسپ بنا دیندی اے۔

لامبڈا لکھت وچ فنکشن دا حوالہ سیدھا اوہنوں متعین کرن والی
علامتی عبارت نال دے سکدے آں، ایس لئی اوہدا وکھرا ناں رکھن دی لوڑ نہیں۔ایہ کہن دی بجائے کہ ``$f$ اوہ فنکشن ہووے جو
$f(x) = x + 3$ نال تعریف ہوندا اے،'' اسی کہہ سکدے آں:
``$f$ فنکشن $\lambd[x][(x + 3)]$ ہووے۔'' ہور لفظاں وچ
$\lambd[x][(x+3)]$ محض اوس فنکشن دا \emph{ناں} اے جیڑھا اپنی
دلیل وچ تین جوڑدا اے۔ایس عبارت وچ $x$ صرف تھاں رکھن والا متغیر
(dummy variable) اے:اوہی فنکشن $\lambd[y][(y + 3)]$ نال وی
ظاہر کیتا جا سکدا اے۔ایہ لکھت دوجے پیرامیٹراں دے ہوندیاں وی کم
کردی اے۔مثلاً فرض کرو $g(x, y)$ دو متغیراں دا فنکشن اے،تے $k$
اک قدرتی عدد اے۔فیر $\lambd[x][g(x,k)]$ اوہ فنکشن اے جیڑھا
ہر~$x$ نوں~$g(x, k)$ ول بھیجدا اے۔

علامتی عبارت توں ایس طرح فنکشن متعین کرن نوں \emph{لامبڈا تجرید
(lambda abstraction)} کہندے نیں۔لامبڈا تجرید دے دوجے پاسے
\emph{اطلاق (application)} اے۔ جے ساڈے کول فنکشن $f$ ہووے
(مثلاً قدرتی عدداں اُتے تعریف شدہ)،تاں اسی اوس نوں 2 ورگی کسے
قدر اُتے \emph{لاگو} کر سکدے آں۔عام لکھت وچ نتیجے لئی~$f(2)$
لکھدے آں۔

لامبڈا تجرید تے اطلاق نوں کٹھے کیتا جاوے تاں کی ہوندا اے؟ جیہڑی دلیل اُتے فنکشن لاگو ہویا اے، اوہنوں تجرید
کیتے متغیر دی تھاں رکھ کے نتیجے والی عبارت سادہ بنائی
جا سکدا اے۔مثلاً
\[
(\lambd[x][(x + 3)])(2)
\]
نوں~$2 + 3$ تک سادہ کیتا جا سکدا اے۔

ایتھوں تک اسی روایتی تصورات لئی صرف نویاں لکھتاں متعارف کروائیاں
نیں۔پر لامبڈا کیلکولس سیٹ نظریاتی زاویے توں ہور بنیادی وکھرا رُخ
لیندا اے۔ایس ڈھانچے وچ:
\begin{enumerate}
\item ہر شے کسی فنکشن نوں ظاہر کردی اے۔
\item فنکشن لامبڈا تجرید نال تعریف کیتے جا سکدے نیں۔
\item ہر شے دا اطلاق ہر ہور شے اُتے ہو سکدا اے۔
\end{enumerate}
مثلاً جے $F$ لامبڈا کیلکولس دا اک ٹرم اے، تاں $F(F)$ نوں
ہمیشہ معنی خیز منیا جاندا اے۔ایس کھلی چھوٹ والے ڈھانچے نوں
\emph{بلا نوع بندی} لامبڈا کیلکولس کہندے نیں؛ایتھے ``بلا نوع بندی''
دا مطلب اے کہ کِنہوں کِس اُتے لاگو کیتا جا سکدا اے،ایس اُتے کوئی
پابندی نہیں۔

\begin{digress}
اک \emph{نوع بندی والا} لامبڈا کیلکولس وی اے، جیڑا بلا نوع بندی
والے روپ دی اہم قسم اے۔بھاویں کئی پاسیاں توں دونے ملدے نیں،پر نوع
بندی والے نوں کلاسیکی سیٹ نظریاتی ڈھانچے نال ملانا بہوں سوکھا اے،
تے اوس دیاں کجھ خاصیتاں کافی وکھریاں نیں۔

لامبڈا کیلکولس اُتے تحقیق نظریاتی کمپیوٹر سائنس تے پروگرامنگ
زباناں دی بناوٹ وچ مرکزی ثابت ہوئی اے۔John McCarthy دی 1950 دی
دہائی وچ بنائی LISP اوہناں مُڈھلیاں زباناں وچوں اک اے جنہاں اُتے ایہناں خیالاں دا اثر پیا۔
\end{digress}

\end{document}
